% Pattern: proof/lead-target. A proof with no title whose body opens by
% naming what it proves in parentheses, as "(of Theorem \ref{b})" or
% "\textit{(Lemma \ref{a})}", proves that item; a title of \thref names it too.
% seen in the gold pilot (set E)
% want lemma-1 statement/lemma 1 labels=a
% want theorem-1 statement/theorem 2 labels=b
% want proof-1 proof - proves=theorem-1
% want proof-2 proof - proves=lemma-1
% want proof-3 proof - proves=theorem-1
\documentclass{amsart}
\newtheorem{theorem}{Theorem}
\newtheorem{lemma}[theorem]{Lemma}
\begin{document}
\begin{lemma}\label{a} Each tree is a graph. \end{lemma}
\begin{theorem}\label{b} Each path is a tree. \end{theorem}
\begin{proof} (of Theorem \ref{b}) A path has no cycle. \end{proof}
\begin{proof} \textit{( Lemma~\ref{a})} A tree is connected. \end{proof}
\begin{proof}[Proof of \thref{b}] Again a path has no cycle. \end{proof}
\end{document}
